\section{Scripture, Mystagogy, Ethics, and Liturgical Song}

Most of Ambrose's surviving treatises began as spoken sermons in the basilica of Milan, delivered on Sundays, feast days, and daily during Lent, then revised from stenographers' notes into books. Reading across his exegetical, sacramental, moral, and liturgical corpus reveals a single pastoral mind: the former magistrate treated every department of Christian learning---cosmology, Gospel commentary, baptismal catechesis, clerical ethics, penitential mercy, and congregational chant---as an instrument for forming an orthodox, praying city.

\subsection{Preaching Creation and the Gospel}

Ambrose's \work{Hexameron} consists of nine homilies on the six days of creation (Genesis 1:1--26), preached during the final week of Lent (probably in 389) and arranged into six books. Like St.~Basil's Greek \work{Hexaemeron}, which he kept close at hand, Ambrose read Genesis aloud to an ordinary congregation, urging them not to lose themselves in speculative astronomy when Scripture points toward eternal life (\work{Hexameron} I.6.20, VI.2.7--8). Yet Ambrose's voice is unmistakably his own---lyrical, concrete, attentive to Italian rivers like the Ticino and Adige, and steeped in liturgical memory. At the close of the fifth day's preaching, as dusk fell and nocturnal birds appeared, he turned to the crowing of the cock before dawn:

\begin{quote}
The cock's song is sweet in the night---not sweet only, but useful---for like a good housemate it wakes the sleeper, warns the anxious, and consoles the traveler, declaring the passage of the night with a tuneful sign. At its singing the robber abandons his ambush; roused by it the morning star rises and lights the sky; at its singing the frightened sailor puts aside his gloom and the storm stirred by evening gales grows calm; at its singing devout affection leaps up to pray and renews the duty of reading; at its singing, last of all, the very Rock of the Church washed away the guilt he had contracted by denial before the cock crowed. (\work{Hexameron} V.24.88, translated from the Latin text of Migne PL~14)
\end{quote}

Every image in that prose paragraph---the sleeper, the morning star, the robber, the sailor, and \latin{ipsa Ecclesiae petra}, ``the very Rock of the Church,'' weeping for his denial---mirrors Ambrose's own Lauds hymn \work{Aeterne rerum conditor} (``Eternal Maker of all things''). He closed that homily with a prayer directly addressing Christ: \latin{Respice nos quoque, Domine Jesu \ldots\ Da, quaeso, lacrymas Petri}---``Look upon us also, Lord Jesus \ldots\ Give, I pray, the tears of Peter'' (\work{Hexameron} V.25.89).

In his ten-book \work{Expositio Evangelii secundum Lucam} (\work{Exposition of the Gospel according to Luke}, c.~386--387), Ambrose gathered homilies into his longest New Testament commentary, moving fluidly among the literal, moral, and mystical senses of Scripture. Commenting on Mary's swift journey into the hill country (Luke 1:39), he distilled the readiness of faith into a single Latin sentence (Schenkl, CSEL 32/4): \latin{nescit tarda molimina sancti spiritus gratia}---``the grace of the Holy Spirit knows no slow endeavors'' (\work{Expositio Lucae} II.19). On Elizabeth's blessing of Mary (Luke 1:45--47), he linked Marian motherhood directly to every believer's interior receptivity: ``If according to the flesh there is one mother of Christ, yet, according to faith, Christ is the fruit of all. For every soul receives the word of God if only he be unspotted and free from sin, and preserves it with unsullied purity'' (\work{Expositio Lucae} II.26, 1843 Oxford \work{Catena Aurea}). And contemplating the Christ Child laid in the manger (Luke 2:7), Ambrose framed the exchange of the Incarnation in words that later echoed through Western Easter liturgy:

\begin{quote}
He was a little child, an infant, that you might be able to become a perfect man; wrapped in swaddling bands, that you might be loosed from the snares of death; laid in a manger, that you might be set upon the altars; on earth, that you might dwell among the stars. \ldots\ I owe more, O Lord Jesus, to Thy sufferings that I am redeemed, than to Thy works that I am created. It were no advantage to be born, had it not advantaged me to be redeemed also. (\work{Expositio Lucae} II.41, opening translated from CSEL 32/4; closing sentences from the 1843 Oxford \work{Catena Aurea})
\end{quote}

\subsection{Mystagogy and the Milanese Altar}

During Lent, Ambrose preached daily on the moral lives of the Patriarchs and the Proverbs, withholding full sacramental explanation until Easter week, after the newly baptized had passed through the font and approached the altar for the first time (\work{De Mysteriis} I.1--2). That reserve---the ancient \latin{disciplina arcani}, or discipline of the secret---served both reverence and pedagogy: ``the light itself of the Mysteries will shed itself with more effect upon those who are expecting they know not what, than if any discourse had come beforehand'' (\work{De Mysteriis} I.2).

In \work{De Mysteriis} (\work{On the Mysteries}, c.~387), Ambrose walked the newly baptized step by step through the rites they had just undergone in the detached Milanese baptistery: the touching of the ears with \latin{Epphatha}, ``Be opened'' (I.3); the renunciation of the devil facing west and the turning toward Christ in the east (II.5--7); the triple immersion in water hallowed by the Cross and the Trinity (III--V); the post-baptismal anointing on the head (VI.29--30); the distinctive Milanese rite of washing the neophytes' feet after baptism, which Ambrose linked to cleansing ``hereditary sins'' inherited from the first man while baptism washes away personal guilt (VI.31--32); the white robes (VII.34); and the sevenfold ``seal of the Spirit'' in Confirmation (VII.42).

When the white-robed neophytes entered the basilica for their first Communion, Ambrose confronted the question every eye instinctively raised:

\begin{quote}
Perhaps you will say, ``I see something else, how is it that you assert that I receive the Body of Christ?'' And this is the point which remains for us to prove. \ldots\ Let us prove that this is not what nature made, but what the blessing consecrated, and the power of blessing is greater than that of nature, because by blessing nature itself is changed. \ldots\ For that sacrament which you receive is made what it is by the word of Christ. \ldots\ Shall not the word of Christ, which was able to make out of nothing that which was not, be able to change things which already are into what they were not? \ldots\ The Lord Jesus Himself proclaims: ``This is My Body.'' Before the blessing of the heavenly words another nature is spoken of, after the consecration the Body is signified. He Himself speaks of His Blood. Before the consecration it has another name, after it is called Blood. And you say, Amen, that is, It is true. (\work{De Mysteriis} IX.50, 52, 54)
\end{quote}

Beside \work{De Mysteriis} stands a longer, more colloquial series of six Easter-week sermons titled \work{De Sacramentis}. The Maurist editors accepted \work{De Sacramentis} as Ambrose's work, while nineteenth-century critics such as Maximilian Ihm rejected it (as de Romestin's Prolegomena notes); twentieth-century patristic scholarship has widely regarded the sermons as unrevised stenographic transcripts of Ambrose's own Easter preaching. What the text of \work{De Sacramentis} IV.5.21--23 and IV.6.27 (Migne PL~16) establishes securely is that late-fourth-century Milan already prayed a Eucharistic canon whose petitions---offering an oblation \latin{ascriptam, ratam, rationabilem, acceptabilem} (``approved, ratified, reasonable, and acceptable''), recounting the institution \latin{Qui pridie quam pateretur}, remembering the Passion, Resurrection, and Ascension, and begging that the offering be borne to the heavenly altar like the gifts of Abel, Abraham, and Melchisedech---stand in direct ancestral continuity with the Roman Canon, without proving that the Milanese and Roman prayers were identical.

\subsection{Clerical Ethics and Penitential Mercy}

Around 391, Ambrose composed his three books \work{De Officiis} (\work{On the Duties of the Clergy}), deliberately modeling its structure---the honorable (\latin{honestum} or \latin{decorum}), the expedient (\latin{utile}), and their unity---on Cicero's Stoic treatise \work{De Officiis}. The imitation was an intentional Christian transfiguration of Roman civic philosophy. Where Cicero measured duty and advantage within the horizon of the earthly republic, Ambrose relocated both into eternal life, replacing Greek and Roman heroes with biblical exemplars---Job, David, Tobit, the Maccabees, and the Christian martyrs---and setting forth the four cardinal virtues of prudence, justice, fortitude, and temperance for the clergy he had ordained.

A similarly pastoral temper governs his two books \work{De Paenitentia} (\work{Concerning Repentance}, c.~384), written against the Novatianist schismatics who styled themselves \latin{katharoi} (``the pure'') and refused ecclesiastical absolution to Christians who fell into grave post-baptismal sin. Ambrose opened Book~I by naming gentleness as the chief virtue of Church rulers: ``he who endeavours to amend the faults of human weakness ought to bear this very weakness on his own shoulders'' (\work{De Paenitentia} I.1.2). Against the Novatianist claim that refusing pardon honored God's majesty, Ambrose answered sharply from John 20:22--23 that Christ gave His priests both binding and loosing together: ``Now no one can repent to good purpose unless he hopes for mercy'' (I.1.4), and ``The Church holds fast its obedience on either side, by both retaining and remitting sin; heresy is on the one side cruel, and on the other disobedient; wishes to bind what it will not loosen, and will not loosen what it has bound'' (I.2.7). Grave sin demanded serious weeping---``greater sins are washed away by greater weeping'' (I.3.10)---yet never despair.

\subsection{The Birth of Western Liturgical Hymnody}

Perhaps Ambrose's most enduring pastoral invention grew out of crisis. During Empress Justina's siege of the Milanese basilicas in 385--386, Ambrose kept his congregation singing antiphonal psalms and newly composed Latin hymns through the night watches. Augustine, then living in Milan, recorded both the innovation and its rapid spread: ``At this time it was instituted that, after the manner of the Eastern Church, hymns and psalms should be sung, lest the people should pine away in the tediousness of sorrow; which custom, retained from then till now, is imitated by many, yea, by almost all of Thy congregations throughout the rest of the world'' (\work{Confessions} IX.7.15).

Ambrose chose a crisp, memorizable meter---eight four-line stanzas of unrhymed iambic dimeter---that packed Nicene Trinitarian confession into lines any artisan or child could sing in the dark. His Homoian opponents mocked the hymns as spells; Ambrose embraced the charge in his \work{Sermon against Auxentius}:

\begin{quote}
They declare also that the people have been led astray by the strains of my hymns. I certainly do not deny it. That is a lofty strain, and there is nothing more powerful than it. For what has more power than the confession of the Trinity which is daily celebrated by the mouth of the whole people? All eagerly vie one with the other in confessing the faith, and know how to praise in verse the Father, Son, and Holy Spirit. So they all have become teachers, who scarcely could be disciples. (\work{Sermon against Auxentius} 34)
\end{quote}

Because later centuries attached Ambrose's name to dozens of liturgical hymns in the same meter, critical scholarship distinguishes a narrow, solidly attested core from the vast ``Ambrosian'' hymnal. Four hymns carry explicit fourth- and fifth-century external authentication (summarized in de Romestin's Prolegomena~\S VI): \work{Aeterne rerum conditor} (for cockcrow/Lauds, mirrored in \work{Hexameron} V.24.88 and cited by Augustine in \work{Retractationes} I.21), \work{Deus creator omnium} (for evening), \work{Iam surgit hora tertia} (for Terce, cited by Augustine), and \work{Intende qui regis Israel} / \work{Veni Redemptor gentium} (for Christmas, cited by Pope Celestine~I in 430 against Nestorianism). When Augustine lay awake at Ostia in 387 mourning his mother Monica, it was the opening stanzas of Ambrose's evening hymn that rose unbidden to soothe his grief:

\begin{quote}
\latin{Deus creator omnium, / Polique rector, vestiens / Diem decora lumine, / Noctem sopora gratia; / Artus solutos ut quies / Reddat laboris usui, / Mentesque fessas allevet, / Luctusque solvat anxios.} --- ``O God, Creator of all things and Ruler of the heavens, clothing the day with fair light and the night with the grace of sleep, that rest may restore loosened limbs to the service of toil, lift weary minds, and unbind anxious griefs.'' (Quoted in St.~Augustine, \work{Confessions} IX.12.32; literal prose rendering)
\end{quote}
